\begin{proof}
\begin{enumerate}[label=\arabic*.]
\item \textbf{The rational connection and the marked paths.}
Put $a=256/3125=-\lambda_*$ and
\[
 D(\lambda)=\operatorname{diag}(1,\lambda^2,\lambda,1,
                              \lambda^{-1},\lambda^{-2}).
\]
For the row-frame equation $Y'=Y\Gamma_6$ of P1.15, the exact
change $\widetilde Y=YD$ gives
\[
 \widetilde\Gamma
 =D^{-1}\Gamma_6D+D^{-1}D'
 =\frac{R_0}{\lambda}+\frac{R_*}{\lambda+a},\qquad
 R_0=\lim_{\lambda\to0}\lambda\widetilde\Gamma,
 \quad R_*=\lim_{\lambda\to-a}(\lambda+a)\widetilde\Gamma.
\]
These formulas specify both rational residue matrices entry by entry
from (1.15.1)--(1.15.2).  Direct rational multiplication verifies the
identity; both complete matrices are also stored in the output named
below.  Their maximum absolute row-sum norms are
\[
 \|R_0\|_\infty=\frac{3125}{256},\qquad
 \|R_*\|_\infty=\frac{14575}{768}.
\]
Let $r=1/32$.  The zero meridian is the polygon with successive
vertices
\[
 1,\ r,\ r+ir,\ -r+ir,\ -r-ir,\ r-ir,\ r,\ 1.
\]
The infinity meridian has successive vertices
\[
 1,\ 1-i,\ -1-i,\ -1+i,\ 1+i,\ 1.
\]
The first runs counterclockwise about $0$ and excludes $-a$; the
second runs clockwise about both finite poles, hence positively
about infinity.  The zero square has radial distance between $r$
and $\sqrt2r<a$, and the infinity square between $1$ and $\sqrt2$,
with $1>a$.  Linear interpolation of these positive radial distances
to $r$ and $1$, respectively, gives based homotopies through the
corresponding punctured regions to their circular meridians; the
zero stem remains the positive segment.  Thus these are the specified
based meridians of P1.15, with matrices $B$ and $A$, respectively.
Since $D(1)=I$, the gauge
keeps their based matrices unchanged.  The open comparison path
from $1$ to $2$ is the positive real segment.

\item \textbf{The Taylor recurrence and its complete error bound.}
Write $M_1=R_0+R_*$ and $M_0=aR_0$.  At a rational complex center
$z$, set
\[
 Q_0=z(z+a),\quad Q_1=2z+a,\quad P_0=zM_1+M_0.
\]
For the fundamental solution $T(t)=\sum_{n\ge0}T_nt^n$ normalized
by $T_0=I$, multiplication of
$T'=T(M_1(z+t)+M_0)/((z+t)(z+a+t))$ gives, with $T_{-1}=0$,
\[
 (n+1)Q_0T_{n+1}
 =T_nP_0+T_{n-1}M_1-nQ_1T_n-(n-1)T_{n-1}.
\]
This recurrence uses only rational connection coefficients and
outward-rounded complex balls.

For each step put
\[
 d_0=\max(|\operatorname{Re}z|,|\operatorname{Im}z|),\quad
 d_*=\max(|\operatorname{Re}(z+a)|,|\operatorname{Im}z|),
 \quad R=\tfrac12\min(d_0,d_*).
\]
The segment is subdivided with rational increments $h$ satisfying
$|h|_1=|\operatorname{Re}h|+|\operatorname{Im}h|\le R/4$.
The whole complex disc $|t|\le R$ avoids both poles, and on it
\[
 \|\widetilde\Gamma(z+t)\|_\infty\le
 M=\frac{\|R_0\|_\infty}{d_0-R}
   +\frac{\|R_*\|_\infty}{d_*-R}.
\]
The integral equation and Gronwall's estimate
\cite[Lemma~2.7, pp.~42--43]{TeschlODE2012}, followed by Cauchy's
coefficient formula
\cite[\S1.9(iii), (1.9.30)--(1.9.31)]{DLMFCauchy2026}, give the
entrywise degree-$N$ tail bound
\[
 E_N=\exp(MR)\frac{(|h|_1/R)^{N+1}}{1-|h|_1/R}.
 \tag{1.16.1}
\]
Adding the real and imaginary error balls of radius $E_N$ to each
entry therefore encloses the exact transition.  All subsequent
matrix products and inverses retain those radii.  The order starts
at $32$ and is increased by $16$ until $E_N<10^{-100}$.
There are $120$ steps for the zero meridian, $66$ for the infinity
meridian and $6$ for the open path, hence $192$ in total; the
largest order is $192$.  Every rational center, increment, radius,
order, norm bound and error ball is saved.

\item \textbf{The normalized intertwiner and the height.}
Let $H_0,H_\infty$ denote the two computed based transition matrices.
The exact equations are
\[
 XH_0=BX,\qquad XH_\infty=AX,\qquad Xe_0=e_0.
\]
P1.15 proves their simultaneous existence and uniqueness.  The
fixed first column leaves thirty unknowns.  We select thirty of
the sixty scalar equations in the other five columns and certify
their coefficient matrix inverse in complex-ball arithmetic.
The exact common solution therefore belongs to the resulting
matrix ball; the remaining equations are also checked to contain
zero.  Their consistency is supplied by P1.15.

Likewise, the real intersection in (1.15.3) is computed from an
invertible four-equation subsystem in
$(\operatorname{Re}v_2,\operatorname{Im}v_2,
  \operatorname{Re}v_1,\operatorname{Im}v_1)$.
P1.8 proves the existence and uniqueness of the full intersection;
the unused equation is checked with its interval radius retained.
At $\lambda=1$ the resulting height ball has radius less than
$4\cdot10^{-91}$.  It is contained in the following wider displayed
interval:
\[
\begin{aligned}
 h_b(1)\in\bigl[&
 -1.03427662241632179949424884537834605259496477790218821132\\
 &\qquad\mathbin{\pm}8.11\cdot10^{-57}\bigr].
\end{aligned}
\]
It lies strictly between the two decimal endpoints in the statement.
Since $5/(3\pi)<1$, we have $h_b(1)+5/(3\pi)<0$.
The cusp asymptotic in P1.8 gives $h_b+c\to+\infty$ along a path
to the cusp.  Continuity on a path from $1$ to that region therefore
supplies a zero for $c=5/(3\pi)$, proving its inadmissibility.

\item \textbf{The regular-root evaluation and the exact comparison.}
The second method uses the regular frames with powers
$(0,6,3,0,-3,-6)$ for $\lambda=u^4$ at zero and
$(0,-8,-4,0,4,8)$ for $\lambda=u^{-5}$ at infinity.
Their transformed connections have denominators
$3125u^4+256$ and $3125+256u^5$.  The regular-root recurrences,
finite deck characters and normalized intertwiner construction
are displayed in the proof of P1.18.  They specify $X_{\rm root}$
and the constant infinity frame $F_\infty$ without integrating a
closed loop.

Let $T_{12}$ be the Fuchsian transition on the positive segment.
The first method gives
\[
 Y_2=XT_{12}D(2)^{-1}.
\]
For the second method put $u=2^{-1/5}$ and
$D_\infty(u)=\operatorname{diag}(1,u^{-8},u^{-4},1,u^4,u^8)$.
If $U_\infty'=U_\infty\Gamma_\infty$ and $U_\infty(0)=I$, then
\[
 Y_{2,\rm root}=F_\infty U_\infty(2^{-1/5})
                         D_\infty(2^{-1/5})^{-1}.
\]
The explicit polynomial numerator of $\Gamma_\infty$ is obtained
by the stated gauge and chain rule from $\Gamma_6$; its coefficient
matrices are saved.  On $|u|\le5/4$, its denominator has modulus
at least $2343.75$, and the row-sum connection bound is $65.6375$.
The degree-$880$ Taylor polynomial and (1.16.1) consequently give
an error less than $5.39\cdot10^{-103}$ at $2^{-1/5}$.
The complete ball comparisons give
\[
 \max_{i,j}|(X-X_{\rm root})_{ij}|<1.37\cdot10^{-73},\qquad
 \max_{i,j}|(Y_2-Y_{2,\rm root})_{ij}|<4.71\cdot10^{-74}.
\]
Thus both period matrices agree to more than thirty decimal digits
under the explicit frame maps above.

\item \textbf{The period planes and the saved certificate.}
For the real parameter $c$ the actual period matrix is
\[
 W_c=(Y_{*,1},\,Y_{*,2},\,Y_{*,3}-(625\pi/96)c\,e_0).
\]
Conjugation of $\det(W_c,\overline{W_c})$ interchanges two blocks
of three columns, so it multiplies the determinant by $(-1)^9=-1$.
Thus the determinant is purely imaginary.  At the two points above
and $c=10$, its positive imaginary parts are
\[
\begin{gathered}
 88279007523.2167451846066607088299709479845\ldots,\\
 3178499608.3637608517583236010752861132441\ldots .
\end{gathered}
\]
The complete stored balls have radii less than $4\cdot10^{-79}$
and $2\cdot10^{-80}$, respectively.  These are evaluations at the
two specified points.  The global admissibility of $c=10$ follows
from the complete minimum argument in P1.18.  The regular-root
calculation there also gives
\[
\begin{split}
 h_b(\infty)&=-1.094540250311402780116268263054945703725533\ldots,\\
 h_b(0)&=-0.931793376356523505430795105796011114990828\ldots .
\end{split}
\]

The rational-gauge, closed-loop and comparison calculation is
\path{code/current_audits/audit_paper_i/closure_07/period_loop_comparison.py};
its complete output is
\path{period_loop_comparison.output.json} in that directory.
The two regular-root inputs are
\path{code/paper_v/output/fresh_frobenius.json} and
\path{fresh_root_connections.json} in the same output directory;
their SHA-256 values are recorded in the comparison output.
The program uses \texttt{python-flint 0.9.0} with $140$ decimal
working digits.  The dependency versions, reproduction command,
exact method and input paths are given in
\path{code/current_audits/audit_paper_i/closure_07/README.md}.
The separate \path{validate_saved_period_comparison.py} and its
JSON output in that directory verify the complete rational path
receipts and all displayed weaker bounds in $980$ exact or ball
checks; the Taylor propagation is supplied by the first program.
\end{enumerate}
\end{proof}
